% 本文件由 LaTeX 排版 Agent 根据有效 translation.json 自动生成。
% author=Bardesanes; work_id=0862; title=诸国法律之书

\chapter{\WorkTitle{诸国法律之书}}

% structure: p0001 source=h1 target=omitted reason=page-title-already-rendered-by-parent

\Signature{Bardesan}\par

数日前，我们去探望我们的兄弟舍玛什格拉姆，巴尔德桑来了，在那里找到了我们。他问候了他的健康，得知他安好之后，便问我们：\Quote{你们刚才在谈论什么？因为我进来时在外面听到了你们的声音。}因为他有一个习惯，每当他发现我们在他来之前谈论什么事情时，他就会问我们：\Quote{你们刚才说什么？}以便与我们谈论此事。\par

\Quote{阿维达在此，}我们对他说道，\Quote{刚才对我们说：\InnerQuote{如果天主是独一的，如你所说，并且他是人的造物主，并且他的旨意是要你们遵行他所命令的，那么他为什么不这样造人，使人不能行恶，而能经常行善呢？因为这样他的旨意就会实现了。}}\par

\Quote{告诉我，我儿阿维达，}巴尔德桑对他说，\Quote{为何你心里会认为万有之神不是独一的？或者他是独一的，却不愿意人们行为公正正直？}\par

\Quote{先生，}阿维达说，\Quote{我问了这些与我同龄的\Emph{弟兄}，好叫\InnerQuote{他们}可以给我一个答复。}\par

\Quote{如果，}巴尔德桑对他说，\Quote{你愿意学习，那么向比你年长的人学习对你更有利；但若是为了教导，你就不需要问他们，而\Emph{相反}应当促使他们问你他们所愿意的。因为教师是被人\InnerQuote{问}问题的，他们自己并不问问题；或者，如果他们曾问问题，那也只是为了引导提问者的心思，使他能恰当地提问，并使他们知道他的愿望是什么。因为人知道如何提问是一件好事。}\par

\Quote{就我而言，}阿维达说，\Quote{我愿意学习，但我先是开始问我这里的弟兄们，因为我不好意思问你。}\par

\Quote{你说得得体，}巴尔德桑说道。\Quote{但你仍要知道，一个恰当地提出问题、愿意被说服、并心平气和地接近真理之道的人，无须感到羞耻；因为他确信，通过我提到的这些事情，他能取悦于他所提问的人。因此，我儿，如果你对这件事有什么见解，请告诉我们大家；如果我们也赞同，我们将表示同意；如果我们不赞同，我们有义务向你说明我们为什么不赞同。但如果你只是希望了解这个主题，而自己对它没有见解，正如一个刚刚加入门徒行列、新近求问的人，我将告诉你\Emph{关于这事}；这样你就不致空手离开我们。此外，如果你对我将要对你说的话感到满意，我们还有其他关于这事的话要对你说；但如果你不满意，我们这一方面也算是陈述了我们的观点，而无任何个人情绪。}\par

\Quote{我也，}阿维达说，\Quote{将非常高兴听到并被说服：因为我不是从别人那里听说这个主题，而是我自己在心里对我的弟兄们说了；他们没有在意说服我，反而说：\InnerQuote{只要相信，你\Emph{就}能够知道一切。}但就我而言，我无法相信，除非我被说服。}\par

\Quote{不仅，}巴尔德桑说，\Quote{阿维达不愿相信，还有许多\Emph{其他人}也因为没有信德在他们里面，甚至不能被说服；但他们总是拆毁又建造，\Emph{如此}他们发现自己对真理一无所知。但尽管如此，既然阿维达不愿相信，看哪！我要对你们这些相信的人讲他问的这件事；\Emph{这样}他也将听到更多有关这事的话。}\par

于是他开始对我们\Emph{如下}讲话：有许多人没有信德，没有从真智慧领受知识。因此，他们不能说话并教导别人，也不容易自己去听。因为他们没有信德的根基可以建造，也没有任何把握可以寄托他们的望德。而且，因为他们连对天主都习惯于怀疑，他们心中也没有对祂的敬畏，这敬畏本身能解救他们脱离一切\Emph{其他}的恐惧：因为心中没有对天主的敬畏的人，是一切\Emph{种类}恐惧的奴隶。因为即使对于那些他们不相信的各种事物，他们也不确定自己不信得正确，而是心意不定，没有确定的信仰，他们思想的滋味在他们\Emph{自己}口中是乏味的；他们总是被恐惧萦绕，情绪激动，鲁莽行事。\par

至于\Person{阿维达}所说：\Quote{天主为何没有如此造我们，使我们不至于犯罪而受审判？}——若人如此受造，他就不属于自己，而是成了那推动他者的工具；显然，那按己意推动\Emph{工具}者，或善或恶地推动它。果真如此，人岂不与别人弹奏的竖琴，或他人驾驶的船无异？在那里，赞美与责难都在演奏者或掌舵者手中，竖琴不知自己被弹奏什么，船也不自知驾驶得好或\Emph{坏}，它们不过是制作来为那拥有\Emph{必要}技能者所用的工具。但天主却仁慈地选择不如此造人；相反，他藉自由将人高举超过许多\Emph{他的}受造物，\Emph{甚至}使他与天使平等。\Emph{请看}太阳、月亮、星辰，以及一切在方面比我们更大的受造物，\Emph{并看}它们如何被剥夺了个体自由，如何全由命令固定于\Emph{其轨道}，以致只能遵行所命，别无其他。因为太阳从未说：我不在指定之时升起；月亮也说：我不变化、不减损、不盈亏；没有星辰说：我不升起也不落下；海也不说：我不承载船只，不守边界；山岳不说：我不留于原位；风不说：我不吹拂；地不说：我不承受并维系其上一切。然而，这一切都是仆人，服从于一个命令：因为它们是永不错误的天主上智的工具。\par

\Emph{但人却不然：}因为若万物皆服侍，谁是被服侍者？若万物皆被服侍，谁是服侍者？\Emph{那样}，便无物相异于他物；然而那单一而无\Emph{部分}差异者，仍是尚未形成的存在。那些注定服侍之物已被置于人的权下：因为人是按\OriginalTerm{厄罗音}{Elohim}的肖像造的。因此，这些事物在\Emph{天主的}仁慈中被赐予人，为暂时服侍他。人也蒙赐由自己意志引导，以致凡他所能行的，若愿意便可为，若不愿意便不为，\Emph{如此}他得以成义或定自己的罪。因为若他受造为不能行恶而招致定罪，同样，他所行的善便不属自己，他也无法因之成义。因为若有人并非自愿行善或行恶，他的成义与定罪便仅系于他所服从的运势。\par

由此你们必明见，天主对人的仁慈何其伟大，自由赐予人的程度远胜于我们所说的一切元素之体，为使人能藉此自由成义，以神圣方式安排行为，并与同样拥有自由意志的天使同为伙伴。因为我们确信，倘若天使也没有自由意志，他们就不会与人的女子交合而犯罪，从高位坠落。同样，那些\Emph{天使}遵行了主的旨意，因着自制而升到更高地位，得以圣化，领受高贵恩赐。因为一切存在之物皆需万有之主；他的恩赐亦无穷尽。\par

然而，\Emph{我所说的}即使那些我称为凭命令存在之物，也并非完全缺乏自由；因此，在末日它们将全受审判。\par

\Quote{但是，}我对他说，\Quote{那些由命令固定\Emph{且由命令规整的}事物，如何受审判呢？}\par

\Quote{并非因它们是固定的，哦，\Person{斐理伯}，}他说，\Quote{那些元素将受审判，而是因它们赋有力量。因为受造之物被形成时，并未被剥夺其自然特性，而\Emph{只是}被削减了力量的完全运用，因彼此混合而\Emph{力量}减弱，并受造物主的力量所制伏；就它们受制伏而言，它们不受审判，但就那\Emph{属于}它们自身\Emph{控制}的而言，则受审判。}\par

\Quote{你所言的这些，} \Person{阿维达}对他说，\Quote{十分美好；但是，看啊！人所接受的诫命是严厉的，他们不能遵行。}\par

\Quote{这话}，\Person{巴德桑}说，\Quote{是一个不愿意行正义之人的话；不仅如此，是一个已经向仇敌屈服顺服之人的话。因为人所受的命令，没有一件是他们不能做的。因为摆在我们面前的诫命只有两条，而且是与自由相称、合乎公平的：一条，我们戒绝一切错误、我们不愿别人对自己做的事；另一条，我们应行正义、我们喜爱并乐于别人也对我们做的事。那么，有哪个人软弱到不能避免偷窃、不能避免说谎、不能避免放荡行为、不能避免仇恨和欺骗？因为看啊！这一切都在人的心智支配之下；不依赖身体的力量，而依赖灵魂的意志。因为即使一个人贫穷、生病、年老、肢体残疾，他也能避免做这一切。而且，正如他能避免做这些事，他也能爱、祝福、说实话，为每一个他所认识的人祈求善事；如果他健康，又能\Emph{工作}，他也能从他所有的中给予；此外，用身体的力量扶助生病衰弱的人——这也是他能做的。}\par

\Quote{那么，究竟是谁不能做那些缺乏信德的人所抱怨的事，我不得而知。就我而言，我认为正是在这些诫命上，人比任何\Emph{其他事}都有更大的能力。因为它们容易，没有任何情况能阻碍它们履行。因为我们没有被命令搬运沉重的石头、木材或任何其他东西——那只有拥有\Emph{极大}体力的人才能做；也没有被命令建造堡垒和创建城市——那只有君王才能做；也没有被命令驾驶船只——那只有水手才有驾驶的技能；也没有被命令丈量和分配土地——那只有\Emph{土地}丈量员才知道如何做；也没有被命令\Emph{从事}任何那些只有少数人拥有而其余人缺乏的技艺。但按照天主的慈善，赐给了我们没有任何严苛的诫命——凡有生命的人都可以乐意去做。因为没有人行正义之事而不感到喜悦，也没有人戒绝坏事而内心不感到欢喜——除非那些不是为此善事而被造、被称为莠子的人。法官若因一个人不能做某事而责备他，难道不是不公正的吗？}\par

\Quote{你说这些行为，啊，\Person{巴德桑}，} \Person{阿维达}对他说，\Quote{是容易做的吗？}\par

\Quote{对于有意志的人，} \Person{巴德桑}说，\Quote{我已经说过，并且\Emph{仍然}说，它们是容易的。因为这种\Emph{我为之辩护的顺从}，是自由心灵和没有反抗其统治者的灵魂的恰当行为。至于身体的行为，有许多事情阻碍它：尤其是年老、疾病和贫穷。}\par

\Quote{或许，} \Person{阿维达}说，\Quote{一个人可能能够戒绝坏事；但至于行善事，有哪个人有能力\Emph{做这事？}}\par

\Quote{行善比戒恶更容易，} \Person{巴德桑}说，\Quote{行善比戒恶更容易。因为善来自人自身，因此他行善时总是欢喜；但恶是仇敌的工作，因此\Emph{正是因为如此，只有当}人被\Emph{某种邪恶情欲}激动，不在正常自然状态时，他才会做坏事。要知道，我儿，赞美祝福朋友是容易的事；但克制自己不辱骂憎恨的人并不容易：尽管如此，仍是可能的。再者，当人行正义时，他的心灵欢喜，良心安宁，他乐意让每个人看见他所做的。但当人行为失当、犯错时，他烦恼激动，充满愤怒狂怒，灵魂和身体都痛苦；当他处于这种\Emph{状态}心神时，他不愿被任何人看见；连那些他原本欢喜、并伴有\Emph{来自他人的}赞美祝福的事，也被他抛在脑后，而那些令他激动不安的事，则\Emph{因伴随\Emph{有意识}的罪疚的诅咒而变得更加痛苦}。}\par

\Quote{然而，或许有人会说，愚昧人做可憎之事时也感到快乐。\Emph{毫无疑问}：但他们快乐不是因为做这些事\Emph{本身}，也不是因为\Emph{因它们}得到称赞，也不是因为他们\Emph{做这些事}怀着美好的希望；快乐本身也不会长久存留。因为在健康状态下的灵魂，怀着美好希望所\Emph{体验}的快乐是一回事；而在病态下的\Emph{灵魂}，怀着不良希望所体验的快乐是另一回事。因为情欲是一回事，而爱是另一回事；友情是一回事，而友善是另一回事；我们应当毫无困难地理解，那种被称为情欲的虚假情感的仿冒品，即使其中有片刻的享受，也与真正的爱截然不同，后者的享受是永恒、不朽坏、不可毁灭的。}\par

\Quote{阿维达}，我对他说，\Quote{还曾这样说：\InnerQuote{人作恶是由于他的\OriginalTerm{本性}{nature}；因为，如果他不是生来就倾向于作恶，他就不会作恶。}}\par

\Quote{如果所有的人，} \Person{巴德桑}说，\Quote{都行为一致，遵循同一种倾向，那么\Emph{显然}就是他们的本性引导他们，他们就没有我一直对你们说的那种自由。但是，为了让你们理解什么是本性、什么是自由，我将继续告知你们。}\par

人的本性是这样的：出生、成长、到达成年、生育子女、衰老、饮食、睡眠、醒来，\Emph{然后}死去。这些事，因为是出于本性，属于所有的人；不仅属于人，也属于一切动物，其中一些也属于树木。因为这就是物理本性的工作，它按照所受的命令制造、产生并调整一切。我说，在动物的行为中也发现本性得以保持。因为狮子吃血肉，是符合它的本性；因此所有的狮子都是食肉的。羊吃草，因此所有的羊都是吃草的。蜜蜂制造蜂蜜来维持生命，因此所有的蜜蜂都是造蜜的。蚂蚁在夏天为自己积蓄，以便在冬天维持生命，因此所有的蚂蚁都这样做。蝎子用它的刺攻击没有伤害它的人，因此所有的蝎子都这样刺。因此所有的动物都保持它们的本性：食肉的不吃草，食草的不吃血肉。\par

相反，人却不是这样被支配的；在涉及他们身体的事上，他们像动物一样保持本性，但在涉及他们心智的事上，他们做自己所选择的事，作为自由、拥有权能、并\Emph{按照天主的肖像受造}的人。因为有些人吃肉，却不碰面包；有些人在\Emph{各种}肉类中作出区别；有些人根本不食用任何动物的肉。有些人娶自己的母亲、姐妹、女儿为妻；有些人根本不与女人交往。有些人自己施行报复，如同狮子和豹子；有些人攻击没有伤害他们的人，如同蝎子；有些人像羊一样被引领，不伤害引领他们的人。有些人以仁慈行事，有些人以正义行事，有些人以邪恶行事。\par

若有人说他们各人都有如此行的本性，应让他确信并非如此。因为有些人\Emph{从前}是放荡和酗酒的人，当善言的劝诫临到他们时，他们变得纯洁和节制，摒弃了肉身的欲望。而有些人\Emph{从前}行为纯洁节制，当他们背离正确的劝诫，胆敢违抗天主和他们导师的命令时，他们从真理之道堕落，变成放荡和狂欢的人。还有些人堕落之后再次悔改，恐惧\Emph{临到并}停留于他们，他们\Emph{重新}转向自己\Emph{先前}所持守的真理。\par

那么，人的本性是什么？看！所有的人在他们行为和目标上彼此不同，只有那些同心同德的人才彼此相像。但那些人——直到此刻仍被他们的欲望诱惑、被他们的怒气支配——决心将自己所做的任何错事归于他们的造主，好让自己显得无可指摘，使造他们的那一位在\Emph{人的}空谈中承担罪责。他们不考虑本性不受法律约束。因为一个人不会因身材高大或矮小、肤色白或黑、眼睛大或小、或任何身体缺陷而受指责；但他若偷窃、说谎、行欺诈、毒害\Emph{他人}、辱骂、或做\Emph{任何}诸如此类的事，则会受指责。\par

\Quote{从此，显然，对于那些不在我们能力之内、而是来自本性的事，我们绝不因此受谴责，也绝不因此而成义；但那些我们在\Emph{运用}个人自由时做的事，若是正确，我们就成义并配得称赞，若是错误，我们就受谴责并担责。}\par

\Quote{还有其他人说，人受命运法令的支配，以至于\Emph{行事}时而邪恶，时而美善。}\par

\Quote{我也知道，哦，\Person{斐理伯}和\Person{巴里雅玛}，}他对我们说，有\Emph{这样的}人：那些被称为加色丁人的，以及其他人喜欢这种精妙学问的人，就像我自己曾经也是。因为我在别处已说过，人的灵魂渴望知道众人所不知道的事，那些人以此为业；并且所有人们犯下的错误，所有他们做得正确的事，以及所有发生在他们身上的事——无论是关于财富还是贫穷，疾病还是健康，身体的残缺——都是通过那被称为七曜的星辰的统治而临到他们的，并且\Emph{事实上}受它们的支配。但还有一些人断言\Emph{相反}这些事——这种技艺是占星家虚妄的发明，命运根本不存在，只是一个虚名；相反，所有的事，无论大小，都掌握在人手中；身体的残缺和过错只是偶然发生并临到他。但另一些人又说，人无论做什么，都是凭自己的意志，在\Emph{运用}赐给他的自由中行事，而临到他的那些过错、残缺和其他不幸之事，是他从天主那里接受的惩罚。\par

然而，依我浅见，此事似乎\Emph{如此}：这三种意见部分可视为真，部分应视为假而弃绝——视为真，因为人按所见现象发言，也因他们见事物临到他们身上，\Emph{仿佛}出于偶然；弃绝为谬，因天主的智慧过于深奥，非他们所能及——那\Emph{智慧}创立世界，造人，设立掌权者，并赋予万物各自相宜的\Emph{优越程度}。我意即：\Emph{此}权能属于天主、天使、权能者、掌权者、元素、人及动物；但我所言这些\Emph{等级}并非于一切事上都拥有此权能（因掌管一切者惟有一位）；而是如我一直所说的，有些事他们有权能，有些事他们无权能：为叫在他们有权能之事上，可见天主的仁慈；在他们无权能之事上，他们知道有一位高于他们。\par

故此，命运如占星家所言，是\Emph{存在}的。再者，万物并非全由我们意志掌控，这由以下情况可见：大多数人都愿富有、统治他人、身体健康、随心所欲支配事物；然而财富只与少数人同在，统治权只在此处彼处一人，身体健康并非人人皆有；\Emph{就连}那些富人也未能完全拥有其财富，掌权者也未能随心所欲支配事物，反而有时事物不顺从\Emph{他们}，不如其愿；富人有时期富有如其所愿，另时期变穷，非其所愿；赤贫者居所不如其愿，在世度日不如其喜，贪求\Emph{许多}事物，却\Emph{只}从他们逃离。许多人得子女，却未养大；另有人养大，却未能保有；还有人保有，子女却成了\Emph{父母的}羞辱与忧愁。有人富足如其所愿，却受疾病之苦，非其所愿；另有人健康如其所愿，却受贫困之苦，非其所愿。有人所愿之物丰富，所不愿者稀少；另有人所不愿者丰富，所愿者稀少。\par

于是此事便\Emph{如此}：财富、尊荣、健康、疾病、子女及其他各种愿望对象，均受命运\Emph{支配}，非我们能力所及；\Emph{相反}，我们因合意之事而喜悦欢欣，对不合意之事则被迫而行；从那些不悦之事发生时，可见那些令我们喜悦之事亦非因我们渴望而发生；而是事物如其所是地发生，有些令我们喜悦，有些则否。\par

\Emph{如此}，我们人类被发现：同受本性管理，各受命运不同支配，又各人凭自由选择。\par

但如今且让我们进而表明：命运并非掌管一切。\Emph{显然不是}：因为所谓命运，本身\Emph{不过是}天主赐予权能者与元素的某种行进秩序；依此行进与秩序，理智下降与灵魂\Emph{结合}时发生变化，灵魂下降与身体\Emph{结合}时也发生变化；此\Emph{秩序}，以命运和\OriginalTerm{生成}{\GreekText{γένεσις}}之名，成为这\Emph{人之组合部分}中发生变化的动因——这组合正被筛净炼净，为使凡因天主的恩宠及善而受益者得益，并正在\Emph{且将继续}得益，直至万终。\par

因此，身体受本性管理，灵魂也分受其经验与感觉；身体在各行为中既不受命运阻碍，亦不受命运助成。因男子十五岁前不能为父，女子十三岁前不能为母。同样，老年也有规律：女子\Emph{此时}不能生育，男子也失去生育的自然能力；而其他动物同样受本性管理，\Emph{甚至}在我所提到的年龄之前，不仅生育后代，而且亦已年老不能生育，正如人之身体年老亦停止繁衍：命运不能在他们身体无生育自然能力时赐予后代。再者，命运不能保全人身体无饮食而存活；即便有饮食，也不能免除死亡：因为这些及其他许多事，专属本性。\par

但\OriginalTerm{本性}{Nature}的时间和方式尽其全部效用之后，\OriginalTerm{命运}{Fate}便来临并出现在它们中间，产生各种效果：时而帮助本性、加强\Emph{它的力量}，时而又削弱和阻挠它。因此，身体的长成和完善来自本性；但身体的疾病和缺陷却来自本性之外，即来自命运。男女的结合以及他们二者纯粹的幸福来自本性；但仇恨和结合的解体，以及\Emph{此外}，所有那些因对交合的\Emph{自然倾向}而人们在淫欲中所行的不洁和纵欲，均来自命运。生育和子女来自本性；而来自命运的是，有时子女畸形，有时遭遗弃，有时夭折。充足的\Emph{营养}供应，足够所有\Emph{受造物}的身体所需，来自本性；而来自命运的是营养匮乏，以及\Emph{随后的}那些身体所受的苦难；同样，从同一个命运又产生暴食和不必要的奢侈。本性规定年长者应为年轻人作审判，智者应为愚者作审判，强壮者应凌驾于软弱者之上，勇敢者应凌驾于怯懦者之上；但命运却导致少年人凌驾于年长者之上，愚者凌驾于智者之上，并且在战争时期软弱者命令强壮者，怯懦者命令勇敢者。\par

你必须清楚地明白，凡在本性偏离其正道之处，其偏离皆因命运而起；\Emph{之所以如此}，是因为那些被称作\OriginalTerm{本命}{Nativity}之变化媒介所依附的首领和治理者，彼此对立。其中有些被称为\OriginalTerm{右方者}{Dexter}，是帮助本性并增强其优势的，每当行运有利于它们，且它们立于那些上升的黄道区域、处于自身宫位之时。反之，那些被称为\OriginalTerm{左方者}{Sinister}则属邪恶，每当轮到它们占据上升位时，它们便与本性作对；它们不仅对人施加伤害，有时也对动物、树木、果实、一年的收成、水源，\Emph{简言之}，对凡在本性之内、受它们治理的一切，都施加伤害。\par

\Quote{因此——\Emph{即}，存在于\OriginalTerm{掌权者}{Potentates}之间的分歧和派别——有些人认为世界\Emph{由这些争斗势力}统治，没有任何\Emph{来自上界}的监督。\Emph{但那}是因为他们不理解，正是这件事——\Emph{我意思是}，\Emph{存在于他们之间的}派别和分歧——以及\Emph{因其行为而产生的}成义和定罪，属于那基于\OriginalTerm{自由}{Freedom}——由\Person{天主}所赐的自由——而被建立的事物的构造，目的是使这些行动者同样，因其自决的力量，得以被成义或被定罪。正如我们看见命运摧毁本性，同样我们也能看见人的自由战胜并摧毁命运本身——然而，并非在一切事上——正如命运本身也并非在一切事上战胜本性。因为这三者，即本性、命运和自由，应当继续存在，直至\Emph{我先前说过的}行运完成，且\Emph{指定的}尺度和\Emph{其演变的}数目得以成就，正如那安排所有受造物的生活方式和结局、以及所有存有和本性的境况者认为合适那样。}\par

\Quote{我深信，} \Person{阿维达}说，\Quote{由您提出的论点，我深信人作恶并非出自其本性，并且并非所有人都受到相同的统治。若您还能进一步证明，那些作恶者并非因命运和\OriginalTerm{定数}{Destiny}而行恶，那么我们就必须相信，人拥有个人的自由，并且按其本性有能力\Emph{既}追随正确的事又避免错误的事，因此也将在末日被公正地审判。}\par

\Quote{您是否，} \Person{巴尔代桑}说，\Quote{由于并非所有人都受到相同的统治这一事实，就确信人作恶并非出自其本性？那么，若我们能向您表明，命运们和掌权者的判决并不同样影响所有人，而是我们在自身内拥有自由，从而能够避免服侍物理的本性和受掌权者控制的影响，您就不可能逃避这样的确信：他们也并非完全出于命运而行恶。}\par

\Quote{请您证明这一点，} \Person{阿维达}说，\Quote{我就信服您，并且无论您命令我做什么，我都愿意去做。}\par

\Quote{您是否，} \Person{巴尔代桑}说，\Quote{读过\Place{巴比伦}的占星家们的书，其中描述了星辰在它们\Emph{不同}的组合中于人的本命时会产生什么效应；以及埃及人的书，其中描述了人可能具有的所有\Emph{各异}的性格？}\par

\Quote{我读过占星术的书，} \Person{阿维达}说，\Quote{但不知道哪些是\Place{巴比伦}人的，哪些是埃及人的。}\par

\Quote{两国的学说，} \Person{巴尔代桑}说，\Quote{是一样的。}\par

\Quote{众所周知是的，} \Person{阿维达}说。\par

\Quote{那么，请听，} \Person{巴尔代桑}说，且请观察：星辰按其命运和宫位所规定的事，并非被在大地所有\Emph{地区}的所有人同样实行。因为人在各国\Emph{为自己}制定了法律，在\Emph{行使}天主所赐予的自由时：因为这一恩赐按其本性就反对那从掌权者发出的命运，掌权者擅自夺取了那未赐予它们的东西。我将开始列举\Emph{这些法律}，尽我所能记起\Emph{它们}，从东方，整个世界的起点开始：——\par

\Emph{塞里斯人的法律}。——塞里斯人有法律禁止杀人、行不洁、拜偶像；在整个塞里卡没有偶像，没有娼妓，也没有杀人者或被杀害者；尽管他们像其他人一样，不分时辰和日子出生。因此，凶悍的火星，无论何时位于天顶，都不能压制塞里斯人的自由，并迫使一个人用铁器杀死同胞；金星与火星合位时，也不能迫使塞里斯人中的任何人与其邻人之妻或任何\Emph{别的}女人交合。然而，那里有富人和穷人、病人和健康人、统治者和被统治者；因为这些事情被交在统治者的权下。\par

\Emph{印度婆罗门人的法律}。——在印度人中，婆罗门人数成千上万，有一条法律完全禁止杀人、敬拜偶像、行不洁、吃肉、饮酒；在这些人民中，这些事情\Emph{从来}不会发生。数千年过去了，这些人一直在他们为自己制定的这条法律之下。\par

\Emph{印度的另一条法律}。——在印度还有另一条法律，在同一地带，盛行于那些不属婆罗门种姓、不接受他们教导的人中，命令他们事奉偶像、行不洁、杀人以及做其他恶事，这些被婆罗门所不齿。在印度的同一地带，还有人有吃人肉的习惯，正如其他国家的人吃兽肉。因此，凶星没有强迫婆罗门行恶和不洁之事；也没有善星能使其他印度人戒绝恶行；也没有那些位于本域、处于有利于仁厚性情的黄道宫中的星辰，能使那些吃人肉的人戒除这种污秽可憎的食物。\par

\Emph{波斯人的法律}。——波斯人又为自己制定法律，准许他们娶姐妹、女儿、孙女为妻；有些人更进一步，甚至娶母亲为妻。这些波斯人的一部分散居海外，\Emph{离开他们的国家}，在\Place{米底亚}、\Place{帕提亚}地区、\Place{埃及}和\Place{弗里吉亚}（他们被称为\Person{玛哥}）\Emph{可以找到}；并且在他们\Emph{可以找到}的所有国家和地带，他们都受这条为他们祖先制定的法律管辖。然而我们不能说，对于所有玛哥和其他波斯人，金星与月亮和土星合位于土星宫中在其分度内，而火星正朝向它们。同样，在帕提亚王国许多地方，男人杀死妻子、兄弟、儿女而不受惩罚；而在罗马人和希腊人中，杀死其中一人者则处以极刑，即最严厉的惩罚。\par

\Emph{革利人的法律}。——在\Place{革利}人中，妇女播种、收割、建造房屋，并做一切劳工的工作；她们不穿彩色衣服，不穿鞋，不涂香料；当她们与陌生男子交合或与家中奴仆亲近时，也没有人指责她们。但这些\Place{革利}妇女的丈夫却穿着彩色衣服，佩戴金饰和宝石，涂上香料。他们这样做并不是因为他们有什么女子气，而是因为为他们制定的法律：事实上，所有男人都喜欢打猎并沉溺于战争。但是我们不能说，对于所有革利妇女，金星位于摩羯宫或宝瓶宫的不利位置；我们也不能说，对于所有革利人，火星和金星位于白羊宫，据说那里出生勇敢而放荡的人。\par

\Emph{巴克特里亚人的法律}。——在\Place{巴克特里亚}人（称为\Place{卡沙尼}人）中，妇女用男人的华丽衣服、大量金饰和贵重宝石装扮自己；奴隶和婢女服侍她们多于服侍她们的丈夫；她们骑马，马匹饰以金鞍和宝石。此外，这些妇女不守节制，与她们的奴仆以及前往该国的陌生男子亲近；她们的丈夫不指责她们，妇女自己也毫不害怕\Emph{受罚}，因为卡沙尼人只\Emph{将}他们的妻子视为情妇。然而我们不能说，对于所有巴克特里亚妇女，金星、火星和木星合位于天中火星之宫，那里出生富有且犯奸淫的妇女，并使她们的丈夫在一切事上顺从。\par

\Emph{拉卡米人、埃德萨人和阿拉伯人的法律}。——在\Place{拉卡米}人、\Place{埃德萨}人和\Place{阿拉伯}人中，不仅犯奸淫的妇女要被处死，就连有奸淫嫌疑的妇女也要受死刑。\par

\Emph{哈特拉的法律}。——在\Place{哈特拉}有一条现行法律：凡偷盗任何细小之物，即使如水一般无价值，都要被石头砸死。\Emph{相反}，在\Place{卡沙尼}人中，如果有人犯了这样的偷窃，他们\Emph{仅仅}朝他脸上吐唾沫。\Emph{同样}，在罗马人中，犯有小偷窃的人被鞭打后打发走。在幼发拉底河对岸，\Emph{再往}东，一个人若被指为小偷或杀人犯，并不太在意；但如果被指为鸡奸者，他会报复甚至将\Emph{指控他的人}杀死。\par

\Emph{法律}.... 在..男孩...给我们，而不... 再有，在整个东方地区，如果有人被\Emph{这样}指控，并被确知\Emph{有罪}，他们的\Emph{亲生}父亲和兄弟就会将他们处死；而且常常连\Emph{埋葬之地}也不让人知道。\par

这些就是东方民族的法律。但是在北方和高卢地区及其邻近地区，那里的男子娶貌美的青年为妻，甚至为此举行宴会；由于他们中间通行的法律，这在他们看来并不算是耻辱或责备。但所有高卢人中被打上这种耻辱印记的人，不可能在出生时水星与金星同在土星宫、位于火星界限内并于西方黄道星座内。因为关于在这种条件下出生的人，经上记载他们被打上不名誉的烙印，如同女人一般。\par

不列颠人的法律。——在不列颠人中，许多男子同娶一个妻子。\par

帕提亚人的法律。——在帕提亚人中，另一方面，一个男子娶许多妻子，而她们全都只忠于他一个人，因为那个国家在那里制定的法律。\par

阿玛宗人的法律。——至于阿玛宗人，她们全体，整个民族，都没有丈夫；而像动物一样，每年春天，她们从自己的领土出发，渡过河流；过河后，她们在一座山上举行盛大节日，那地区的男子们前来与她们共度十四天，与她们交往，她们因此受孕，然后再次返回本国；分娩时，所生的男婴她们就丢弃，女婴则抚养长大。显然，按照自然规律，既然她们都在同一月受孕，那么她们也在同一月分娩，或稍早或稍晚；正如我们所听说的，她们全都强壮好战；但没有任何星辰能够帮助那些出生的男婴免遭丢弃。\par

\Quote{\WorkTitle{占星家之书}。——在占星家的书中记载，当水星与金星同在水星宫时，产生画家、雕刻家和银行家；但当它们同在金星宫时，产生香料商、舞者、歌手和诗人。\Emph{然而}，在整个塔伊特人和撒拉森人地区，在上利比亚和毛里塔尼亚人中间，在海洋入口处的诺马德人地区，在外日耳曼尼亚，在上萨尔马提亚，在西班牙，在蓬图斯以北的所有国家，在整个阿兰人地区，在阿尔巴尼亚人中间，在扎齐人中间，在杜罗河以外的布鲁萨，既见不到雕刻家、画家，也见不到香料商、银行家或诗人；相反，水星和金星的这项法令被阻止\Emph{影响}整个世界的圆周。在整个米底亚，所有人死后，\Emph{甚至}当生命尚存时，就被丢给狗吃，狗吃掉了整个米底亚的死者。然而我们不能说所有米底亚人出生时月亮与火星同在巨蟹座、在日间地下：因为经上记载，被狗吃的人是这样出生的。印度人死后，全部被火焚烧，他们的许多妻子也活活同焚。但我们不能说所有被焚烧的印度妇女在出生时火星和太阳同在狮子座、在夜间地下，如同那些被火烧死的人那样。所有日耳曼人，除了战死的人，都死于绞刑。但不可能所有日耳曼人出生时月亮和\OriginalTerm{霍拉}{Hora}都处于火星和土星之间。事实上，在所有国家，每天每时，人们出生时的星位各不相同，人的法律胜过星辰的法令，他们受自己的习俗支配。命运不能强迫塞里斯人违背意愿去谋杀，也不能强迫婆罗门吃肉；它不能阻止波斯人娶自己的女儿和姐妹，不能阻止印度人被焚烧，不能阻止米底亚人被狗吃，不能阻止帕提亚人娶许多妻子，不能阻止不列颠人中许多男子同娶一个妻子，不能阻止厄德撒人修持贞洁，不能阻止希腊人从事体操，……不能阻止罗马人不断夺取其他土地，不能阻止高卢的\Emph{男人}互相结婚；它也不能强迫阿玛宗人养育男婴；他的星位也不能强迫整个世界圆周内的任何人从事缪斯艺术；但是，正如我\Emph{已经}说过的，在每个国家和每个民族，所有的人都按自己所选择的方式利用他们本性的自由，并由于他们所披戴的身体，有时按自己的意愿，有时违反自己的意愿，为命运和本性服务。因为在每个国家和每个民族，都有富人和穷人、统治者和臣民、健康者和患病者——各人按命运和\Emph{他的}星位对他的影响。}\par

\Quote{关于这些事，神父巴尔代桑，}我对他说，\Quote{您说服了我们，我们知道这是真的。但是您知道吗，占星家们说大地分为七部分，称为区域；而那七颗\Emph{星}在那些部分之上各有权柄，每颗\Emph{掌管一个}；在每一部分中，其统治者的意志占主导；这被称为它的\Emph{法律}？}\par

\Quote{首先，你要知道，我儿斐理伯，}他对我说，星象学家们捏造这种说法，作为\Emph{促进}谬误的伎俩。因为，虽然大地分为七部分，但在每一部分中，都可以发现许多彼此不同的法律。因为世上并非只有七种法律，按照七星的数目；也不是十二种，按照黄道十二宫的数目；也不是三十六种，按照旬星的数目。而是，当你从一国到另一国、从一地区到另一地区、从一区到另一区、在每个\Emph{人的}居所，都可以看到许多\Emph{种类的}法律，彼此不同。因为你还记得我对你说过的话——在一个地带，即印度人的地带，有许多人不吃动物的肉，还有一些人\Emph{甚至}吃人肉。再者，我告诉过你，\Emph{在谈及}波斯人和贤士时，不仅是在波斯的地带，他们\Emph{娶}自己的女儿和姐妹为妻，而且在他们所到的每个国家，他们都遵守祖先的法律，并保守他们传授给他们的那种\Emph{教导}中所包含的神秘技艺。再者，请记住我告诉过你，有许多民族散布在全世界的整个圆形区域，他们并不局限于任何一个地带，而是居住在四面八方、所有地带，并且他们没有墨丘利和维纳斯\Emph{被认为是}在彼此会合时所给予的技艺。然而，如果法律受地带支配，这是不可能的；但显然它们并不如此：因为我所提及的\Emph{那些人}与许多\Emph{其他}人在生活习惯上毫无共通之处。\par

\Emph{再者}，你们想想，有多少智者废除了他们认为制定得不好的本国法律？又有多少法律因必要而被搁置？还有多少国王，当他们占领了不属于他们的国家时，废除了既定的法律，并制定了\Emph{其他}他们想要的法律？每当这些事情发生时，没有一颗星辰能够保全法律。这里有一个例子就近在眼前，你们可以\Emph{亲眼}看到：就在昨天，罗马人占领了阿拉伯，他们废除了\Emph{那里}先前存在的所有法律，特别是他们所行的割礼。事实上，自主的人\Emph{有时}被迫服从他人强加给他的法律，而他人\Emph{反过来}又获得了为所欲为的权力。\par

但让我向你们提及一个事实，它比\Emph{其他}任何事物都更能说服愚昧和缺乏信心的人。所有通过梅瑟领受法律的犹太人，都在第八天为他们的男婴施行割礼，而不等待\Emph{合适的}星辰到来，也不畏惧他们\Emph{所居住}国家的法律。掌管地带的星辰也无力强制他们；无论他们在厄东、阿拉伯、希腊、波斯、北方还是南方，他们都遵守祖先为他们制定的这条法律。显然，他们的行为并非出自本命星：因为不可能所有犹太人都在第八天行割礼时，火星均\Quote{当升}，以至于刀锋加于他们身上，流血。此外，他们所有人，无论在哪里，都禁止向偶像致敬。每七天中的一天，他们和他们的儿女停止所有工作，停止一切建筑、旅行、买卖；他们在安息日不杀生、不生火、不施行审判；在他们中间找不到任何人被命运强迫，或在安息日去诉讼并赢得官司，或去诉讼而输掉，或拆毁、建造，或做所有没有领受这条法律的人所做的任何事情。他们还有其他事情，在这些事情上，他们在\Emph{安息日}不像其余人类那样行事，尽管在这一天他们同样生育、降生、生病和死亡：因为这些事不属人力范围。\par

在叙利亚和厄德撒，人们曾为了敬礼塔拉塔而自阉；但当阿布加尔王成为信徒后，他下令凡行此事者砍断其手，从那天直到如今，在厄德撒境内无人再行此事。\par

我们基督徒的新民族，基督在祂来临时栽种在每个国家和每个地区，我们又该说什么呢？因为，看！无论我们在哪里，我们都以基督的同一名字——基督徒——被称呼。在每周的第一天，我们聚集在一起，在诵读日，我们禁绝\Emph{进食}。在高卢的弟兄们不\Emph{娶}男性为妻；在帕提亚的弟兄们不娶两个妻子；在犹太的弟兄们不自行割礼；我们那些在格利人中的姐妹们不与外人交往；在波斯的弟兄们不娶自己的女儿\Emph{为妻}；在玛待的弟兄们不抛弃死者，不活埋他们，也不将他们喂狗；在厄德撒的弟兄们不杀死犯不洁之罪的妻子或姐妹，而是离开她们，将她们交托给天主的审判；在哈特拉的弟兄们不将窃贼用石头\Emph{打死}；但是，无论他们在哪里，无论他们在什么地方\Emph{被发现}，\Emph{各国}的法律都不妨碍他们服从他们的\Emph{君主}基督的法律；\Emph{天界}执政者的命运也不能强迫他们使用他们认为不洁的东西。\par

反之，疾病与健康、富有与贫穷，这些不属于他们自由范围的事，无论他们在哪里都会遭遇。因为尽管人的自由并不受七曜的强制影响，或者即使有时受影响，也能抵挡施加于它的影响，然而，\Emph{反之}，\Emph{同}一个人，从外部来看，无法即刻摆脱掌权者对他的命令：因为他是一个奴隶，处于隶属状态。因为，如果我们能做成一切事，我们自己就是一切；如果我们没有任何能力做事，我们就成了他人的工具。\par

\Quote{但是，当天主愿意\Emph{它们}时，一切皆有可能，\Emph{并且它们可以毫无阻碍地发生}：因为没有任何事物能阻挡那伟大而神圣的旨意。因为，即使是那些自以为\Emph{成功地}抵挡它的人，也不是凭力量抵挡，而是凭邪恶和错误。这种情况可能持续一小段时间，因为他善待并容忍一切存在之物，让他们保持原样，凭自己的意愿受管理，然而他们却受已完成的作为和为他们帮助而制定的安排所约束。因为所建立的\Emph{这个}有序的万物构成和治理，以及彼此的混合，抑制了\Emph{这些}存在的暴力，使他们既不能对\Emph{彼此}施以完全的伤害，也不能完全受伤害，就像他们在世界受造之前那样。还有一个时刻即将来临，那时他们里面仍存留的\Emph{施加}伤害的倾向，将通过\Emph{给予他们}的另一种交往中的教导而终结。在那新世界的建立之时，一切邪恶的动荡都将止息，一切悖逆都将终结，愚昧人将被说服，一切缺陷将被弥补，将有宁静与和平，作为万有之主的恩赐。}\par

\Emph{此}处结束《\WorkTitle{诸国律法书}》。\par

因此，\Person{巴尔代桑}，一位年长的人，以其对事件的知识而闻名，在他所著的一部作品中，论及天上光体之间的同时周期，这样写道：——\par

土星的两个周期，60年；\par

木星的五个周期，60年；\par

火星的四十个周期，60年；\par

太阳的六十个周期，60年；\par

金星的七十二个周期，60年；\par

水星的一百五十个周期，60年；\par

月亮的七百二十个周期，60年。\par

他说，这是它们所有的一个同时周期；即它们一个\Emph{这样的}同时周期的时间。因此\Emph{可见}，\Emph{完成}一百个这样的同时周期将\Emph{需要}六千年。如下：——\par

土星的二百个周期，六千年；\par

木星的五百个周期，六千年；\par

火星的四千个周期，六千年；\par

太阳的六千个周期，六千年；\par

金星的七千二百个周期，六千年；\par

水星的一万二千个周期，六千年；\par

月亮的七万二千个周期，六千年。\par

\Person{巴尔代桑}如此计算这些，是想要表明这个世界只存续六千年。\par

\SourceRecord{Translated by B.P. Pratten.；Ante-Nicene Fathers；Vol. 8.；Edited by Alexander Roberts, James Donaldson, and A. Cleveland Coxe.；Buffalo, NY: Christian Literature Publishing Co.,；1886.；<http://www.newadvent.org/fathers/0862.htm>.}
